\chapter{Ergänzende Implementierungs- und Prüfdetails}
\label{app:repro}
\section{Implementierungsauszüge}
Die folgenden Auszüge gehören zu den in Kapitel~5 erläuterten Implementierungsschritten.
Die vollständigen Notebooks und die nachvollziehbaren Prüfskripte liegen in den digitalen
Artefakten. Kürzungen innerhalb der Ausschnitte sind keine lauffähigen Ersatzprogramme.

\par\medskip\noindent\begin{minipage}{\dimexpr\linewidth-1em\relax}
\begin{lstlisting}[language=Python,caption={Zentraler Trainingsaufbau des URL-TAPT -- gekürzter Originalausschnitt},label={lst:url_tapt}]
# KI-Unterstuetzung: siehe KI-Erklaerung der Arbeit.
def train_url_tapt(batch_size):
    seed_all(MASTER_SEED)
    urls = read_ssl_urls()  # labels are never read here

    enc = url_tok(
        urls,
        truncation=True,
        max_length=URL_MAX_LEN,
        padding=False
    )
    model = AutoModelForMaskedLM.from_pretrained(
        URL_BASE_SRC, token=HF_TOKEN, **url_kwargs
    ).to(DEVICE)

    coll = DataCollatorForLanguageModeling(
        tokenizer=url_tok,
        mlm_probability=URL_MLM_PROB
    )
    opt = torch.optim.AdamW(
        model.parameters(),
        lr=URL_TAPT_LR,
        weight_decay=URL_TAPT_WEIGHT_DECAY
    )

    total = len(dl) * URL_TAPT_EPOCHS
    sched = get_linear_schedule_with_warmup(
        opt, int(URL_TAPT_WARMUP_FRAC * total), total
    )
\end{lstlisting}
\end{minipage}\par\medskip

\par\medskip\noindent\begin{minipage}{\dimexpr\linewidth-1em\relax}
\begin{lstlisting}[language=Python,caption={Bestimmung der Development-basierten Kaskadenschwellen},label={lst:cascade_thresholds}]
# KI-Unterstuetzung: siehe KI-Erklaerung der Arbeit.
url_high = thr_fpr(
    ET["url"][yt == 0],
    .001
)

pos = np.sort(ET["url"][yt == 1])
k = max(1, int(math.floor(.01 * len(pos))))
url_low = float(
    np.nextafter(pos[k - 1], -np.inf)
)
\end{lstlisting}
\end{minipage}\par\medskip

\par\medskip\noindent\begin{minipage}{\dimexpr\linewidth-1em\relax}
\begin{lstlisting}[language=Python,caption={Routingentscheidung der URL-first-Kaskade -- gekürzter zusammenhängender Originalauszug},label={lst:cascade_route}]
# KI-Unterstuetzung: siehe KI-Erklaerung der Arbeit.
hi = us >= url_high
lo = us <= url_low
esc = ~(hi | lo)

out[hi] = 1e9
out[lo] = -1e9
\end{lstlisting}
\end{minipage}\par\medskip

\par\medskip\noindent\begin{minipage}{\dimexpr\linewidth-1em\relax}
\begin{lstlisting}[language=Python,caption={Deterministische Rangbildung für die Datenpartitionierung},label={lst:stable_rank}]
# KI-Unterstuetzung: siehe KI-Erklaerung der Arbeit.
def stable_rank(series, salt):
    return np.fromiter(
        (
            int(
                hashlib.sha256(
                    (salt + "|" + x).encode()
                ).hexdigest()[:16],
                16
            )
            for x in series.astype(str)
        ),
        dtype=np.uint64,
        count=len(series)
    )
\end{lstlisting}
\end{minipage}\par\medskip

\par\medskip\noindent\begin{minipage}{\dimexpr\linewidth-1em\relax}
\begin{lstlisting}[language=Python,caption={Aufbereitung von HTML-Text und strukturellen Webseiteninformationen -- gekürzter Originalausschnitt},label={lst:process_row}]
# KI-Unterstuetzung: siehe KI-Erklaerung der Arbeit.
def process_row(row, include_label=True):
    s, _ = clipped_html(row.get("html"))
    no_comments = COMMENT_RE.sub(" ", s)

    toks = []
    for slash, tag in TAG_RE.findall(no_comments):
        toks.append(("/" if slash else "") + normalize_tag(tag))
        if len(toks) >= STRUCT_MAX_TOKENS:
            break
    struct = " ".join(toks)

    txt = SCRIPT_STYLE_RE.sub(" ", no_comments)
    txt = TAG_STRIP_RE.sub(" ", txt)
    txt = html_std.unescape(txt)
    txt = WS_RE.sub(" ", txt).strip()[:TEXT_MAX_CHARS]

    tags, parents, depths, attrs, children, _ = parse_dom(s)
    out = {
        "text": txt,
        "html_struct": struct,
        "template_hash": hashlib.sha256(
            struct.encode("utf-8", "ignore")
        ).hexdigest(),
        "dom_tag": tags,
        "dom_parent_idx": parents,
        "dom_depth": depths,
        "dom_attr_count": attrs,
        "dom_child_count": children,
    }
    return out
\end{lstlisting}
\end{minipage}\par\medskip

\par\medskip\noindent\begin{minipage}{\dimexpr\linewidth-1em\relax}
\begin{lstlisting}[language=Python,caption={Physische Entfernung der Klassenlabels aus dem SSL-Pool},label={lst:remove_label}]
# KI-Unterstuetzung: siehe KI-Erklaerung der Arbeit.
for role, outdir in ROLE_DIRS.items():
    x = df[df._role.eq(role)].copy()
    if not len(x):
        continue

    if role == "SSL_POOL":
        x["ssl_rank"] = x["_ssl_rank"].astype(int)
        x = x.drop(columns=["label"], errors="ignore")

    x = x.drop(columns=["_role", "_ssl_rank"], errors="ignore")
    x.to_parquet(
        outdir / f"part-{counters[role]:05d}.parquet",
        index=False,
        compression="zstd"
    )
    counters[role] += 1
\end{lstlisting}
\end{minipage}\par\medskip

\par\medskip\noindent\begin{minipage}{\dimexpr\linewidth-1em\relax}
\begin{lstlisting}[language=Python,caption={Partielle Freigabe der Transformerencoder},label={lst:freeze_last}]
# KI-Unterstuetzung: siehe KI-Erklaerung der Arbeit.
def freeze_last(model, n):
    for p in model.parameters():
        p.requires_grad = False

    layers = getattr(
        getattr(model, "encoder", None),
        "layer",
        None
    )
    if layers is None:
        layers = getattr(
            getattr(model, "transformer", None),
            "layer",
            None
        )
    if layers is None:
        raise RuntimeError(type(model))

    for layer in layers[-n:]:
        for p in layer.parameters():
            p.requires_grad = True

    if getattr(model, "pooler", None) is not None:
        for p in model.pooler.parameters():
            p.requires_grad = True
\end{lstlisting}
\end{minipage}\par\medskip

\par\medskip\noindent\begin{minipage}{\dimexpr\linewidth-1em\relax}
\begin{lstlisting}[language=Python,caption={GCN-basierte Repräsentation des statischen DOM -- gekürzter Originalausschnitt},label={lst:dom_encoder}]
# KI-Unterstuetzung: siehe KI-Erklaerung der Arbeit.
class DOMEncoder(nn.Module):
    def __init__(self):
        super().__init__()
        self.tag = nn.Embedding(len(DOM_VOCAB), 96, padding_idx=PAD)
        self.depth = nn.Embedding(32, 16)
        self.attr = nn.Embedding(32, 8)
        self.child = nn.Embedding(32, 8)
        self.inp = nn.Linear(128, DOM_DIM)
        self.layers = nn.ModuleList(
            [GCNLayer(DOM_DIM) for _ in range(DOM_LAYERS)]
        )
        self.att = nn.Linear(DOM_DIM, 1)

    def nodes(self, g):
        x = self.inp(torch.cat([
            self.tag(g["tag"]), self.depth(g["depth"]),
            self.attr(g["attr"]), self.child(g["child"])
        ], 1))

        for layer in self.layers:
            x = layer(x, g["edge"])
        return x
\end{lstlisting}
\end{minipage}\par\medskip

\par\medskip\noindent\begin{minipage}{\dimexpr\linewidth-1em\relax}
\begin{lstlisting}[language=Python,caption={Rekonstruktionsverlust des DOM-SSL -- gekürzter Originalausschnitt aus dem Trainingsloop},label={lst:dom_ssl}]
# KI-Unterstuetzung: siehe KI-Erklaerung der Arbeit.
logits = m(g)
mask = g["mask"]
if mask.any():
    loss = F.cross_entropy(
        logits[mask],
        g["orig"][mask]
    )
    opt.zero_grad(set_to_none=True)
    loss.backward()
    torch.nn.utils.clip_grad_norm_(
        m.parameters(), 1
    )
    opt.step()
\end{lstlisting}
\end{minipage}\par\medskip

\par\medskip\noindent\begin{minipage}{\dimexpr\linewidth-1em\relax}
\begin{lstlisting}[language=Python,caption={Kern der multimodalen DUAL-/TRI-Fusion -- gekürzter Architekturausschnitt},label={lst:deep_fusion}]
# KI-Unterstuetzung: siehe KI-Erklaerung der Arbeit.
class DeepFusion(nn.Module):
    def __init__(self, use_dom):
        super().__init__()
        self.gate = nn.Sequential(
            nn.Linear(PROJ_DIM, 64),
            nn.GELU(),
            nn.Linear(64, 1)
        )

        nm = 3 if use_dom else 2
        self.head = nn.Sequential(
            nn.Linear(PROJ_DIM * (nm + 1), 512),
            nn.GELU(),
            nn.Dropout(.2),
            nn.Linear(512, 128),
            nn.GELU(),
            nn.Linear(128, 1)
        )

    def forward(self, u, t, g=None):
        zs = self.encode(u, t, g)
        st = torch.stack(zs, 1)
        gw = torch.softmax(self.gate(st).squeeze(-1), 1)
        weighted = (st * gw[:, :, None]).sum(1)
        main = self.head(
            torch.cat(zs + [weighted], 1)
        ).squeeze(-1)
        return main, gw
\end{lstlisting}
\end{minipage}\par\medskip

\section{Reproduktionsumfang}
Die digitale Fassung enthält die Ergebnisdateien der Primäranalysen, getrennte ergänzende
Versuchsblöcke und die beim abschließenden Review verwendeten Prüfskripte.
\nolinkurl{empirical_audit.py} rekonstruiert die aggregierten Kennzahlen und statistischen
Entscheidungen; \nolinkurl{raw_score_audit.py} prüft Vorhersagescores gegen wiedergewonnene
Testlabels, reproduziert die stratifizierten Bootstrap-Intervalle und trainiert die
externen Linear-Probes auf den gespeicherten Repräsentationen neu.

Die Neutrainings verwenden unveränderte Auswahlhashes und Modellvorgaben. Sie ersetzen
keine Wiederholung des selbstüberwachten Encodertrainings. Einzelne Rangmetriken können
bei einer erneuten Optimierung numerisch abweichen; die Differenzen zu den historischen
Exporten werden im beigefügten Prüfbericht ausgewiesen. Die Latenz- und Durchsatzprüfung
kann nur die exportierten Messaggregate, deren Rechenbeziehungen und den Messcode
nachvollziehen, da die einzelnen Zeitmessungen nicht exportiert wurden.
